\renewcommand{\bibname}{Quellen und Anschlüsse im Projekt}
\begin{thebibliography}{99}
\addcontentsline{toc}{chapter}{Quellen und Anschlüsse im Projekt}
\small
\interlinepenalty=10000

\bibitem{projectbracketing}
Martin Kunze: \emph{Die Grundlagen der Mathematik}, Band~28,
Ergänzung \emph{Klammerungsunabhängigkeit}.
Projektquellen: \path{tex/b28/bracketing/}.
Die allgemeine Assoziativität verbindet diesen Leitfaden mit der
Entwicklung endlicher Wörter und Klammerungsbäume in Band~27.

\bibitem{projectpower}
Martin Kunze: \emph{Die Grundlagen der Mathematik}, Band~28,
Abschnitte über Potenzhalbgruppen und deren Homomorphismen;
Ergänzung \emph{Halbgruppenbeweis der Potenzhalbgruppe}.
Projektquellen: \path{Bd. 28 - Halbgruppen.tex},
\path{tex/b28/power/}, \path{tex/b28/power-isomorphisms/}.

\bibitem{harju}
Tero Harju: \emph{Lecture Notes on Semigroups. Short Course}.
Universität Turku, 2010.
Kapitel~1: Grundbegriffe, Erzeugnisse und Homomorphismen;
Kapitel~2: Kongruenzen und Ideale;
Kapitel~3: freie Halbgruppen und Präsentationen;
Kapitel~4: allgemeine Green-Relationen und der Gruppensatz
als vollständige Alternative für jede $\mathcal H$-Klasse.
\url{https://users.utu.fi/harju/semigroups/Semigroups2010.pdf}.

\bibitem{gould}
Victoria Gould: \emph{Semigroup Theory}.
Vorlesungsskript, University of York.
Insbesondere Satz~2.6 (Transformationsdarstellung),
Proposition~2.15 (monogene Halbgruppen) und
die allgemeinen Aussagen zu Green-Relationen in den Abschnitten~5--6.
\url{https://www-users.york.ac.uk/~varg1/SemigroupTheory.pdf}.

\bibitem{howie}
John M. Howie: \emph{Fundamentals of Semigroup Theory}.
London Mathematical Society Monographs, New Series~12,
Clarendon Press, Oxford, 1995.
Grundbegriffe, Kongruenzen und Green-Relationen;
Satz~3.1.6 zur Hauptfaktorstruktur.
\url{https://doi.org/10.1093/oso/9780198511946.001.0001}.

\bibitem{green1951}
J. A. Green: \emph{On the Structure of Semigroups}.
Annals of Mathematics, Second Series, 54 (1951), 163--172.
Historische Originalquelle der Green-Relationen und ihrer Strukturtheorie.
\url{https://doi.org/10.2307/1969317}.

\bibitem{bajri}
Sanaa Ahmed E Bajri: \emph{On Representations of Semigroups}.
Dissertation, University of York, 2020.
Abschnitt~4.1, insbesondere Theorem~4.1.1 und Lemma~4.1.2,
zur Struktur der Hauptfaktoren, nach Howie.
\url{https://etheses.whiterose.ac.uk/id/eprint/27806/1/Bajri_202041702_CorrectedThesisClean.pdf}.

\bibitem{burris}
Stanley Burris und H. P. Sankappanavar:
\emph{A Course in Universal Algebra}.
Graduate Texts in Mathematics~78, Springer, 1981;
frei zugängliche überarbeitete Fassung.
Kapitel~I, Abschnitt~5, sowie Kapitel~II, Sätze~3.2--3.3 und~5.5
(S.~30--31 und~37--38), zu algebraischen Unterstruktur- und
Kongruenzverbänden; Abschnitt~8 zu subdirekten Produkten;
Satz~11.9 (S.~75--76) zum HSP-Satz;
Satz~14.19 (S.~95) und Übung~11 zu Abschnitt~14 (S.~98)
zur Vollständigkeit der Gleichungslogik und zu endlichen Ersetzungsketten.
Die hier verwendeten Aussagen sind auf die Signatur einer einzigen
assoziativen zweistelligen Operation spezialisiert.
\url{https://www.math.uwaterloo.ca/~snburris/htdocs/UALG/univ-algebra2012.pdf}.

\bibitem{harju96}
Tero Harju: \emph{Lecture Notes on Semigroups}.
Universität Turku, Vorlesungsskript, 1996, kleine Korrekturen 2010.
Abschnitt~2.1, Lemma~2.3, zur syntaktischen Kongruenz
einer beliebigen Teilmenge einer Halbgruppe.
\url{https://users.utu.fi/harju/semigroups/semigroups96.pdf}.

\bibitem{brookes2024}
Matthew Brookes, James East, Craig Miller, James D. Mitchell und Nik Ruškuc:
\emph{Heights of one- and two-sided congruence lattices of semigroups}.
Pacific Journal of Mathematics 333 (2024), 17--57.
Abschnitt~5.4: allgemeine Grundlagen zu Schützenbergergruppen.
\url{https://doi.org/10.2140/pjm.2024.333.17}.

\bibitem{koussas2021}
Timothy Koussas: \emph{The Classification of Dualisable Aperiodic Semigroups}.
Dissertation, La Trobe University, 2021.
Abschnitt~2.1, Propositionen~2.1--2.2, S.~12--13:
allgemeine Korrespondenz zwischen Idealen und unteren Mengen
der $\mathcal J$-Klassen.
\url{https://opal.latrobe.edu.au/ndownloader/files/31212412}.

\bibitem{nybergbrodda2023}
Carl-Fredrik Nyberg-Brodda:
\emph{On the Word Problem for Free Products of Semigroups and Monoids}.
Journal of Algebra 622 (2023), 721--741.
Abschnitt~1.3, insbesondere Lemma~1.1: allgemeine Normalform
des freien Halbgruppenprodukts.
\url{https://arxiv.org/abs/2112.10665}.

\bibitem{munn1968}
W. D. Munn: \emph{Regular $\omega$-semigroups}.
Glasgow Mathematical Journal 9 (1968), 46--66.
Abschnitt~1.5, S.~48: Bruck-Erweiterung eines beliebigen Monoids
und Einbettung jeder Halbgruppe; dort nach R.~H. Bruck,
\emph{A Survey of Binary Systems} (1958), Satz~8.3.
\url{https://doi.org/10.1017/S0017089500000288}.

\bibitem{bogdanovic2011}
Stojan M. Bogdanović, Miroslav D. Ćirić und Žarko Lj. Popović:
\emph{Semilattice Decompositions of Semigroups}.
Faculty of Economics, University of Niš, 2011.
Sätze~4.1 und~4.2, S.~126--127, zur feinsten
Halbverbandszerlegung und zur Unzerlegbarkeit ihrer Komponenten.
\url{https://www1.pmf.ni.ac.rs/pmf/vesti/biblioteka/knjige/Semilattice_Decompositions_of_Semigroups.pdf}.

\bibitem{isbell1966}
John R. Isbell: \emph{Epimorphisms and Dominions}.
In: \emph{Proceedings of the Conference on Categorical Algebra,
La Jolla 1965}, Springer, 1966, 232--246.
Satz~2.3: Zigzag-Satz für beliebige Halbgruppen.
\url{https://doi.org/10.1007/978-3-642-99902-4_9}.

\bibitem{alam2019}
Noor Alam, Peter M. Higgins und Noor Mohammad Khan:
\emph{Epimorphisms, dominions and H-commutative semigroups}.
Autorenversion, 2019.
Abschnitt~1, insbesondere Result~1.4: allgemeine Dominions
und Isbells Zigzag-Satz, vor den Resultaten zu Spezialklassen.
\url{https://arxiv.org/abs/1908.01813}.

\bibitem{higgins1995}
Peter M. Higgins: \emph{Embedding semigroups in semibands of depth two}.
Semigroup Forum 51 (1995), 395--396.
Allgemeiner Einbettungssatz in eine Halbgruppe, deren Elemente Produkte
zweier idempotenter Elemente sind.
\url{https://doi.org/10.1007/BF02573647}.

\bibitem{oliveira2014}
Luís Oliveira: \emph{On the question of embedding a semigroup into a semiband}.
Semigroup Forum 88 (2014), 715--731.
Autorenfassung von 2012 unter dem Titel
\emph{On the question of embedding a semigroup into an idempotent
generated one}: Einleitung, S.~2, und Abschnitt~2, S.~4,
zum allgemeinen Einbettungssatz von Higgins.
\url{https://arxiv.org/abs/1205.3454};
\url{https://doi.org/10.1007/s00233-013-9564-x}.

\bibitem{howie1981}
John M. Howie: \emph{A class of bisimple, idempotent-generated
congruence-free semigroups}.
Proceedings of the Royal Society of Edinburgh, Section A,
88 (1981), 169--184.
Synopsis, S.~169: universelle Einbettungen in solche Halbgruppen
für beliebig große reguläre unendliche Kardinalzahlen.
Die hier formulierte Existenz einer Erweiterung jeder Halbgruppe
wird im üblichen Rahmen mit Auswahlaxiom verwendet.
\url{https://doi.org/10.1017/S0308210500017376}.

\bibitem{birkhoff1935}
Garrett Birkhoff: \emph{On the Structure of Abstract Algebras}.
Proceedings of the Cambridge Philosophical Society 31 (1935), 433--454.
Abschnitt~10, Theorem~10, S.~441; moderne Formulierung bei
Burris--Sankappanavar, Satz~II.11.9.
Im Leitfaden auf Halbgruppen spezialisiert und mit der vereinbarten
Zulassung leerer Halbgruppen formuliert.
\url{https://doi.org/10.1017/S0305004100013463}.
\bibitem{mitsch1994}
Heinz Mitsch: \emph{Semigroups and their natural order}.
Mathematica Slovaca 44 (1994), 445--462.
Einleitung, S.~445--446: natürliche partielle Ordnung beliebiger
Halbgruppen und äquivalente Definitionen; Satz~2.6 (P.~R. Jones),
S.~449--450, zur Ordnungsisomorphie der Hauptordnungsideale
innerhalb einer $\mathcal D$-Klasse. Abschnitt~3 behandelt
die nicht allgemein gegebene Multiplikationsverträglichkeit.
\url{https://dml.cz/dmlcz/132918}.
\bibitem{howie1968}
John M. Howie: \emph{Commutative semigroup amalgams}.
Journal of the Australian Mathematical Society 8 (1968), 609--630.
Einleitung, S.~609, zum Isbell-Satz über Dominions und zwei Kopien
derselben Halbgruppe; Abschnitt~1, S.~610--612, zum universellen
amalgamierten Produkt. Diese Aussagen gelten für beliebige,
auch nichtkommutative Halbgruppen. Abschnitt~2 behandelt das mögliche
Scheitern der Injektivität bei allgemeinen Amalgamen.
\url{https://doi.org/10.1017/S1446788700006273}.
\end{thebibliography}
